\begin{solution}{104}
\begin{subsolution}
以凳面中心 $O$ 为坐标原点，以过 $O$ 点向上的竖直线为 $y$ 轴。茶杯（包括茶水在内）的质心位置 $y_{CM}$ 为
\begin{equation}
y_{\mathrm{CM}} = \frac{\frac{4m}{5}\frac{H}{2} + \rho(\pi r^{2} h)\frac{h}{2}}{m + \rho(\pi r^{2} h)}
\label{eq:1.s104}
\end{equation}
式中 $h$ 是杯中茶水的高度. 令 $\lambda = \rho \pi r^{2}$ ，\eqref{eq:1.s104} 式即
\[
y_{\mathrm{CM}} = \frac{\frac{4m}{5}\frac{H}{2} + \lambda \frac{h^{2}}{2}}{m + \lambda h}
\]
事实上，考虑在茶杯中茶水的水平面从杯底逐渐缓慢上升的过程，茶水杯整体的质心先是逐渐降低，然后再逐渐上升。为了使茶水杯盛尽可能多的茶水并保持足够稳定，茶水杯整体的质心应尽可能接近凳面，处于最低点的位置 $y_{\mathrm{CM}}(h = h_{\max}) = (y_{\mathrm{CM}})_{\min}$ ，故有
\begin{equation}
\left.\frac{\mathrm{d} y_{\mathrm{CM}}}{\mathrm{d} h}\right|_{h = h_{\max}} = 0
\label{eq:2.s104}
\end{equation}
由\eqref{eq:1.s104}和\eqref{eq:2.s104}式得
\begin{equation}
h_{\max} = \frac{m}{\lambda}\left[-1 \pm \sqrt{1 + \frac{4}{5}\frac{H\lambda}{m}}\right]
\label{eq:3.s104}
\end{equation}
舍弃负值（不合题意），杯中茶水的最佳高度为
\begin{equation}
\begin{array}{r l}
h_{\max} & = \frac{m}{\lambda}\left(\sqrt{1 + \frac{4}{5}\frac{\lambda H}{m}} - 1\right) \\
& = \frac{m}{\rho\pi r^{2}}\left(\sqrt{1 + \frac{4\pi}{5}\frac{\rho r^{2} H}{m}} - 1\right)
\end{array}
\label{eq:4.s104}
\end{equation}

[另一解法:
事实上，考虑在茶杯中的茶水的水平面从杯底逐渐缓慢上升的过程，茶水杯整体的质心先是逐渐降低，然后再逐渐上升. 为了使茶水杯盛尽可能多的茶水并保持足够稳定，茶水杯整体的质心应尽可能接近凳面，处于最低点的位置 $y_{\mathrm{CM}}(h = h_{\max}) = (y_{\mathrm{CM}})_{\min}$ . \eqref{eq:1.s104}式即
\begin{equation}
y_{\mathrm{CM}} = \frac{\frac{2mH}{5} + \frac{m^{2}}{2\lambda}}{m + \lambda h} + \frac{m + \lambda h}{2\lambda} - \frac{m}{\lambda}
\label{eq:5.s104}
\end{equation}
它满足代数不等式
\begin{equation}
y_{\mathrm{CM}} \geq 2\sqrt{\frac{\frac{2mH}{5} + \frac{m^{2}}{2\lambda}}{m + \lambda h} \frac{m + \lambda h}{2\lambda}} - \frac{m}{\lambda} = 2\sqrt{\frac{\frac{2mH}{5} + \frac{m^{2}}{2\lambda}}{2\lambda}} - \frac{m}{\lambda} = (y_{\mathrm{CM}})_{\min}
\label{eq:6.s104}
\end{equation}
式中当 $h$ 满足
\begin{equation}
\frac{m + \lambda h}{2\lambda} = \frac{\frac{2mH}{5} + \frac{m^{2}}{2\lambda}}{m + \lambda h}
\label{eq:7.s104}
\end{equation}
即
\begin{equation}
\begin{array}{r l}
h = h_{\max} & = \frac{m}{\lambda}\left(\sqrt{1 + \frac{4}{5}\frac{\lambda H}{m}} - 1\right) \\
& = \frac{m}{\rho\pi r^{2}}\left(\sqrt{1 + \frac{4\pi}{5}\frac{\rho r^{2} H}{m}} - 1\right)
\end{array}
\label{eq:8.s104}
\end{equation}
时取等号，$h_{\max}$ 即杯中茶水的最佳高度.]
\end{subsolution}

\begin{subsolution}
记凳子质量为 $M$. 该茶水杯的底面边缘刚好滑移到与圆凳的边缘内切于 $D$ 点静止后，坐标系及三脚圆凳的受力分析分别如解题图a和解题图b所示，其中 $N_{A}$、$N_{B}$ 和 $N_{C}$ 分别是车内底板对凳脚 $A$、$B$ 和 $C$ 的支持力．由几何关系有
\begin{equation}
\overline{O' A} = \overline{O' B} = \frac{a/2}{\cos 30^{\circ}} = \frac{a}{\sqrt{3}}
\label{eq:9.s104}
\end{equation}
\begin{equation}
\overline{O' O''} = \frac{a}{2}\tan 30^{\circ} = \frac{a}{2\sqrt{3}}
\label{eq:10.s104}
\end{equation}
三脚圆凳处于力平衡状态，竖直方向合力为零
\begin{equation}
N_{A} + N_{B} + N_{C} - Q - P = 0
\label{eq:11.s104}
\end{equation}
\begin{equation}
N_{A} = N_{C}
\label{eq:12.s104}
\end{equation}
式中，$Q = mg + \lambda h_{\max}g$，$P = Mg$．由\eqref{eq:11.s104}和\eqref{eq:12.s104}式得
\[
2N_{A} + N_{B} - Q - P = 0
\]
以 $x$ 轴为转动轴，力矩平衡
\begin{equation}
Q(R - r) + N_{B}\overline{O' B} - (N_{A} + N_{C})\overline{O' O''} = 0
\label{eq:13.s104}
\end{equation}
记 $R' = R - r$，由\eqref{eq:9.s104}、\eqref{eq:10.s104}和\eqref{eq:13.s104}式有
\[
Q R' + N_{B}\frac{a}{\sqrt{3}} - 2N_{A}\frac{a}{2\sqrt{3}} = 0
\]

[另一解法:
以 $AC$ 为转动轴，由力矩平衡条件有
\begin{equation}
Q\left(R' - \overline{O' O''}\right) + N_{B}\left(\overline{O' B} + \overline{O' O''}\right) - P\overline{O' O''} = 0
\label{eq:14.s104}
\end{equation}
即
\[
Q\left(R' - \frac{a}{2\sqrt{3}}\right) + N_{B}\left(\frac{a}{\sqrt{3}} + \frac{a}{2\sqrt{3}}\right) - P\frac{a}{2\sqrt{3}} = 0
\]
联立以上各式得
\begin{equation}
\begin{array}{r l}
N_{B} & = \frac{(a - 2\sqrt{3}R')Q + aP}{3a} \\
& = \left[\left(1 - 2\sqrt{3}\frac{R - r}{a}\right)\sqrt{1 + \frac{4\pi}{5}\frac{\rho r^{2} H}{m}}\, m + M\right]\frac{g}{3}
\end{array}
\label{eq:15.s104}
\end{equation}
\begin{equation}
\begin{array}{r l}
N_{A} = N_{C} & = \frac{(a + \sqrt{3}R')Q + aP}{3a} \\
& = \left[\left(1 + \sqrt{3}\frac{R - r}{a}\right)\sqrt{1 + \frac{4\pi}{5}\frac{\rho r^{2} H}{m}}\, m + M\right]\frac{g}{3}
\end{array}
\label{eq:16.s104}
\end{equation}
值得指出的是，解\eqref{eq:15.s104}满足 $N_{B}>0$，即
\[
M > \frac{2\sqrt{3}(R - r) - a}{a}\sqrt{1 + \frac{4\pi}{5}\frac{\rho r^{2} H}{m}}\, m
\]
这已由题给条件保证了.]

由\eqref{eq:11.s104}、\eqref{eq:15.s104}和\eqref{eq:16.s104}式得，此时旅行车内底板对各凳脚的支持力相对于茶水杯滑移前的改变为
\begin{equation}
\begin{array}{r l}
\Delta N_{B} & = \frac{(a - 2\sqrt{3}R')Q + aP}{3a} - \frac{P + Q}{3} \\
& = - \frac{2\sqrt{3}}{3}\frac{R - r}{a}\sqrt{1 + \frac{4\pi}{5}\frac{\rho r^{2} H}{m}}\, m g
\end{array}
\label{eq:17.s104}
\end{equation}
\begin{equation}
\begin{array}{r l}
\Delta N_{A} = \Delta N_{C} & = \frac{(a + \sqrt{3}R')Q + aP}{3a} - \frac{P + Q}{3} \\
& = \frac{\sqrt{3}}{3}\frac{R - r}{a}\sqrt{1 + \frac{4\pi}{5}\frac{\rho r^{2} H}{m}}\, m g
\end{array}
\label{eq:18.s104}
\end{equation}
\end{subsolution}
\end{solution}